\chapter{Evaluation}
\label{chap:evaluation}

The evaluation has two parts, a technical evaluation of pitch detection 
and score following, and a user evaluation.

\section{Objectives}

\begin{itemize}
	\item Measure how accurately each pitch detection algorithm identifies notes.
	\item Find the configuration (algorithm, window, FFT size) that works best.
	\item Check that the cursor follows the musician correctly.
	\item Assess the usability of the applicatin with real users.
\end{itemize}

\section{Technical evaluation}

\subsection{Test environment}
\label{subsec:test-env}

We built a browser-based test bench to compare algorithms and configurations 
under the same conditions. It runs the same WebAssembly module as the application.
Each test is a JSON file with an input, an expected result, and optionally its 
own configuration. The bench can run two configurations side by side on the same 
tests.

\todo[inline]{Add screenshot of the test bench.}

\subsection{Test types}

\begin{itemize}
	\item \textbf{Pitch detection:} synthetic sine waves at known frequencies, from a 
	single note (A4, 440\,Hz) to all 120 notes from C0 to B9.
	\item \textbf{Audio recognition:} real instrument recordings with a known sequence 
	of notes. \todo[inline]{add some more details (instruments use, etc)}
	\item \textbf{Score tracking:} checks that the cursor ends at the expected position 
	of the score, within a tolerance. \todo[inline]{add some more details}
\end{itemize}

The score-following logic is also covered by unit tests in Rust. They check that the 
cursor moves to each played note, stays still when the input is silent, ignores notes 
that are not in the score, skips rests and stops at the end of the score.

\subsection{Configurations}

\todo[inline]{Add table comparing configurations, link to algorithms section, 
maybe a screenshot of test bench}

\subsection{Metrics}

\todo[inline]{table explaining the metrics used}

\subsection{Results}

\todo[inline]{explain test bench results, add screenshots of test bench}

\section{User evaluation}

\subsection{Participants}

[Number, instrument, level, age range]

\subsection{Evaluation tasks}

[What each participant does (load a score, play, ...)]

\subsection{Procedure}

[Where, how long, devices used, ...]

THe questionnaire is in Appendix~\ref{ap:cuestionario}

\subsection{Results}

\section{Discussion}

[Findings and limitations of the evaluation]
